\documentclass[10pt,a4paper]{amsart}
\usepackage[margin=19mm]{geometry}
\usepackage{amsmath,amssymb,mathtools}
\usepackage[scr=rsfs,cal=euler]{mathalfa}
\usepackage{xfrac}
\usepackage[unicode,hidelinks,bookmarks=false]{hyperref}
\newcommand{\ie}{i.e.\ }
\newcommand{\cf}{cf.\ }
\newcommand{\resp}{respectively\ }
\newcommand{\Subsection}{\subsection}
\providecommand{\nobreakdash}{\nobreak\hbox{-}}
\setlength{\emergencystretch}{2em}
\multlinegap0pt
\usepackage{bm}
\usepackage{bbm}
\usepackage{dsfont}
\usepackage{xspace}
\usepackage{cancel}
\usepackage{mathtools}
\usepackage{amsthm}
\makeatletter
\@ifundefined{theorem}{\newtheorem{theorem}{Theorem}}{}
\@ifundefined{corollary}{\newtheorem{corollary}[theorem]{Corollary}}{}
\@ifundefined{proposition}{\newtheorem{proposition}[theorem]{Proposition}}{}
\makeatother
\title[Infinite Moments of Class Groups for Solvable Fields with a ]{Infinite Moments of Class Groups for Solvable Fields with a Normal Abelian Subgroup}
\author{Weitong Wang}
\date{Source: arXiv:2606.01831v1 (CC BY 4.0)}
\begin{document}
\maketitle

\noindent\emph{Source and licence.} Infinite Moments of Class Groups for Solvable Fields with a Normal Abelian Subgroup. Version arXiv:2606.01831v1 (CC BY 4.0), distributed under the Creative Commons Attribution 4.0 International licence, as recorded in the arXiv licence metadata for this exact version and shown on the abstract page https://arxiv.org/abs/2606.01831v1. Full rights notice: licenses/arxiv-2606.01831-CC-BY-4.0.txt.\par\smallskip

\noindent\textit{Editorial scope.} Counted unit: the Proposition whose source label is \texttt{prop: 1-1 correspondence for G-structured algebras} in the article (source0.tex lines 692--699, source label \texttt{prop: 1-1 correspondence for G-structured algebras}), delivered together with the article's own complete original proof of it (source0.tex lines 700--745, one \texttt{proof} environment of 46 lines). No mathematics is written, repaired or re-derived: statement and proof are the article's own text line for line. Required context retained uncounted: thm: structure of Galois algebras (lines 263--268); prop: isomorphisms of G-algebras (lines 632--636); cor: Galois G-extensions (lines 673--680). Every definition, display and label of those lines is retained unchanged. Omitted from this package and not redistributed: 1--262, 269--631, 637--672, 681--691, 746--1869. Nothing inside a retained excerpt is omitted or altered: every retained body is delivered exactly as the article's lines are written, so no omission receipt of any kind accompanies this package. Citations: the retained excerpts carry no citation command at all, so no bibliography entry of the article is transcribed and no reference is redistributed. Reference adaptation: none was needed and none was made; every reference inside the delivered unit resolves inside it, so no unresolved reference or citation is produced. Printed slips kept literal: no printed word, symbol, hypothesis or number of the counted statement or of its proof was altered. Layout and class: the article's own class is \texttt{amsart} with packages including amsmath, amssymb, amscd, amsthm, amsfonts, enumerate, tikz, tikz-cd, verbatim, xspace, tensor, graphicx, appendix, color; the wrapper is the lane's pinned \texttt{amsart} editorial wrapper at 10pt on A4, which restates the theorem-like environments the retained text needs and adds the article's own macro definitions that the retained text needs (none), with extra wrapper packages (bm, bbm, dsfont, xspace, cancel, mathtools). The delivered numbering is the unit's own. Assets: no retained span contains a figure environment, a graphics-inclusion command, a PDF-page inclusion, a table or a TikZ picture; no image file is copied into this package, the package declares no asset dependency and no image file is redistributed with it. Licence: version arXiv:2606.01831v1 of this article is distributed under the Creative Commons Attribution 4.0 International licence, as recorded in the arXiv licence metadata for this exact version and shown on the abstract page https://arxiv.org/abs/2606.01831v1. A full rights notice is delivered at licenses/arxiv-2606.01831-CC-BY-4.0.txt. Characters: every retained excerpt is pure ASCII, so the delivered engine, XeLaTeX with its default Latin Modern text font, reports no missing character.
	\begin{theorem}\label{thm: structure of Galois algebras}
		Let $G$ be a finite group, $F$ be a field such that $\operatorname{char}(F)\nmid\lvert G\rvert$, and $(A/F,\varphi_A)$ be a Galois $(G,F)$-algebra.
		For each primitive (central) idempotent $e\in A$, let $G_e:=\operatorname{Stab}_G(e)$ be its stabilizer.
		The component $eA$ is a Galois $G_e$-extension over $F$ with the group isomorphism $\varphi_{eA}:G_e\to\operatorname{Gal}(eA/F)$ defined by $g\mapsto\varphi_A(g)\vert_{eA}$.
		In particular, $A\cong\operatorname{Ind}^G_{G_e} eA$ as Galois $(G,F)$-algebras, where $\operatorname{Ind}$ means the induced algebra.
	\end{theorem}

	\begin{proposition}\label{prop: isomorphisms of G-algebras}
		Let $(A,\varphi_A)$ and $(B,\varphi_B)$ be two Galois $(G,F)$-algebras.
		They are isomorphic as Galois $(G,F)$-algebras if and only if there exists primitive idempotents $e\in A$ and $e'\in B$ such that the Galois $F$-algebras $(eA,G_e)$ and $(e'B,G_{e'})$ are isomorphic by a conjugation from $G$.
		That is, there exists some $g\in G$ such that the pair $(f:eA\to e'B,c_g:G_e\to G_{e'})$ is an isomorphism, where $c_g(h)=ghg^{-1}$ is the conjugation by $g$.
	\end{proposition}

	\begin{corollary}\label{cor: Galois G-extensions}
		Let $K/F$ be a fixed Galois field extension such that $\operatorname{Gal}(K/F)\cong G$.
		Given two group isomorphisms $\varphi_1,\varphi_2:G\to\operatorname{Gal}(K/F)$, the two Galois algebras $(K,\varphi_1)$ and $(K,\varphi_2)$ are isomorphic if and only if there exists some $g\in G$ such that
		\begin{equation*}
			\varphi_2=\varphi_1\circ c_g,
		\end{equation*}
		where $c_g(h)=ghg^{-1}$ is the conjugation by $g$.
	\end{corollary}

	\begin{proposition}\label{prop: 1-1 correspondence for G-structured algebras}
		Fix a finite group $G$ and a field $F$.
		There is a one-to-one correspondence between the following two sets
		\begin{equation*}
			\{\text{Galois }(G,F)\text{-algebras up to isomorphism}\}\leftrightarrow\operatorname{Hom}(G_F,G)/\sim
		\end{equation*}
		where $G_F$ is the absolute Galois group of $F$, and $\sim$ means the equivalence relation induced by the conjugate of $G$, that is, $\rho,\chi:G_F\to G$ are equivalent if there is some $g\in G$ such that $\chi=g\rho g^{-1}$.
	\end{proposition}

	\begin{proof}
		Let $\bar{F}$ be the algebraic closure of $F$, and $\rho:G_F\to G$ be a continuous homomorphism with image $H$.
		Then $K:=\bar{F}^{\ker\rho}$ is a Galois $H$-extension over $F$, where the map $\varphi_K:H\to\operatorname{Aut}_F(K)$ is induced by $\rho$.
		And $A:=\operatorname{Ind}^G_H K$ is naturally a Galois $(G,F)$-algebra.
		
		If $\chi:G_F\to G$ is another continuous homomorphism such that there exists some $g\in G$ for all $\sigma\in G_F$ we have $\chi(\sigma)=g\rho(\sigma)g^{-1}$, then $K=\bar{F}^{\ker\chi}$ is a Galois $gHg^{-1}$-extension over $F$.
		This implies that $(K,H)$ and $(K,gHg^{-1})$ are isomorphic as Galois $F$-algebras by $(1_K,c_g)$.
		Let $B:=\operatorname{Ind}^G_{gHg^{-1}}K$.
		By the Proposition~\ref{prop: isomorphisms of G-algebras}, we see that $A$ and $B$ are isomorphic as Galois $(G,F)$-algebras.
		
		Therefore, we have a well-defined map
		\begin{equation*}
			\Phi:\operatorname{Hom}(G_F,G)/\sim\to\{\text{Galois }(G,F)\text{-algebras}\}/\cong,\quad\varphi\mapsto\operatorname{Ind}^G_{\varphi(G_F)}\bar{F}^{\ker\varphi}.
		\end{equation*} 
		Let $A$ be a Galois $(G,F)$-algebra.
		By the Theorem~\ref{thm: structure of Galois algebras}, if $e$ is a primitive idempotent, then $eA$ is a Galois $(G_e,F)$-field extension.
		As an algebraic extension over $F$, we could choose an embedding $eA\hookrightarrow\bar{F}$.
		Then we have a surjective continuous homomorphism $\rho_{eA}:G_F\to\operatorname{Aut}_F(eA)$,
		and the isomorphism $\varphi:G_e\to\operatorname{Aut}_F(eA)$ induces a continuous homomorphism $\rho_A:=\varphi_A^{-1}\circ\rho_{eA}:G_F\to G$ with image $G_e$.
		By the Corollary~\ref{cor: Galois G-extensions}, we see that different embeddings $eA\hookrightarrow\bar{F}$ induces conjugate surjective homomorphisms $G_F\to\operatorname{Aut}_F(eA)$.
		Therefore the map $A\mapsto\rho_A$ is defined up to conjugation.
		That is, the algebra $A$ corresponds to a class of homomorphisms.
		
		If $B$ is another Galois $(G,F)$-algebra that is isomorphic to $A$, then by the Proposition~\ref{prop: isomorphisms of G-algebras}, we may assume without loss of generality that $e$, resp. $e'$, is a primitive idempotent of $A$, resp. of $B$, such that $G_e=G_{e'}$ and $f:(eA,\varphi_{A})\to(e'B,\varphi_{B})$ is an isomorphism of Galois $(G_e,F)$-algebras.
		Given any embedding $\sigma:e'B\hookrightarrow\bar{F}$, we see that $\sigma\circ f:eA\hookrightarrow\bar{F}$ is an embedding such that the induced maps $\rho_A=\varphi_A^{-1}\circ\rho_{eA}$ and $\rho_B=\varphi_B^{-1}\circ\rho_{eB}$ coincide.
		This implies that $A$ and $B$ correspond to the same class of continuous homomorphisms.
		So, we have a well-defined map
		\begin{equation*}
			\Psi:\{\text{Galois }(G,F)\text{-algebras}\}/\cong\to\operatorname{Hom}(G_F,G)/\sim,\quad(A,\varphi_A)\mapsto\rho_A.
		\end{equation*}
		Then it suffices to check that $\Phi$ and $\Psi$ are inverse to each other.
		Let $\rho:G_F\to G$ be a continuous homomorphism with image $H$ and kernel $G_K$.
		Up to the equivalence relations, we have $\Phi(\rho)=\operatorname{Ind}^G_H K$, denoted by $A:=\Phi(\rho)$.
		The Galois $(G,F)$-algebra $A$ has a component $K$ which is a Galois $H$-extension over $F$.
		It is already a subfield of $\bar{F}$ by its construction, so $\varphi_K:H\to\operatorname{Aut}_F(K)$ induces the map $\rho_A:G_F\to G$ by composing with the restriction of Galois actions $G_F\to\operatorname{Aut}_F(K)$.
		But this is exactly the original map $\rho$, for the map $\varphi_K:H\to\operatorname{Aut}_F(K)$ is induced by $\rho$.
		That is, the action of $H$ on the field extension $K/F$ is given by $\rho:G_F\to G$.
		So, we see that $\Psi\circ\Phi$ is the identity map.
		
		Now let $(A,\varphi_A)$ be a Galois $(G,F)$-algebra, $e$ be a primitive idempotent, and $\sigma:eA\hookrightarrow\bar{F}$ be a choice of embedding with the induced surjective map $\rho_{eA}:G_F\to\operatorname{Aut}_F(eA)$.
		By definition, $\rho_A=\Psi(A)=\varphi_A^{-1}\circ\rho_{eA}$.
		Clearly, $K:=\bar{F}^{\ker\rho_A}=\sigma(eA)$, and $K$ is a Galois $(G_e,F)$-field by the map $\rho_A:G_F\to G$.
		So, we see that $\sigma:eA\to K$ is an isomorphism of Galois $G_e$-field extensions over $F$, and $\Phi(\rho_A)=\operatorname{Ind}^G_{G_e}K$ is isomorphic to $A$ as Galois $(G,F)$-algebras by the Proposition~\ref{prop: isomorphisms of G-algebras}.
		This implies that $\Phi\circ\Psi$ is also the identity (of the opposite direction).
		And we are done.
	\end{proof}

\end{document}
